\documentclass[12pt]{article}
\usepackage[utf8]{inputenc}
\usepackage{geometry}
\geometry{letterpaper, margin=1in}
\usepackage{amsmath}
\usepackage{enumitem}

\title{CS 412 Assignment 2 Report}
\author{Donggyu Park (donggyu5)}
\date{\today}

\begin{document}

\maketitle

\section{Part 1}

\subsection{Part 1 Code (35 points)}

Code is submitted separately via homework2\_q1.py to Gradescope.

\subsection{Short Answers (15 points)}
% Answer in 2 to 3 sentences each.

\begin{enumerate}[label=(\alph*)]
    \item \textbf{Why does bagging primarily reduce variance, while boosting can reduce bias?}
    \\% TODO: Write your answer here (2-3 sentences)
    Bagging typically uses base learners such as deep decision trees, which have high variance. 
    By averaging the predictions from the high variance base learners, bagging can effectively smooth out
    the noise (or high variance) of the individual learners. As a result, bagging reduces overall variance.

    \item \textbf{Why are decision stumps commonly used as base learners in boosting?}
    \\% TODO: Write your answer here (2-3 sentences)
    In Boosting, decision stumps (depth = 1) act as weak learners that have high bias and low variance.
    Boosting builds an ensemble model by sequentially training these weak learners.
    Since ensenble learning builds a strong learner by sequentially training multiple weak learners based on previous errors,
    using simple models like decision stumps allows boosting to effectively reduce bias while controlling variance.

    \item \textbf{What happens to AdaBoost if the base learner performs worse than random guessing? Briefly explain why.}
    \\% TODO: Write your answer here (2-3 sentences)
    In AdaBoost, if the base learner error rate is worse than random guessing $\epsilon > 1 - 1/K$, the algorithm will stop the training.
    If a weak learner performs worse than random guessing, it means that the contribution of that learner will make the overall model perform worse.
    Also, with a high error rate, the model voting weight, $\alpha$, will be negative, making the ensemble learning (more weight on difficult samples) impossible.
    
\end{enumerate}

\vspace{1em}
\hrule
\vspace{1em}

\section{Part 2}

\subsection{Part 2 Code (35 points)}
Code is submitted separately via homework2\_q2.py to Gradescope.

\subsection{Data Challenges: Imbalance and Missing Values (15 points)}

\begin{enumerate}[label=(\alph*)]
    \item \textbf{Class Imbalance (Report Task)}
    % TODO:
        \begin{enumerate}[label=(\arabic*)]
        \item \textbf{Why accuracy is a poor metric for evaluating performance in this scenario.}
        \\ If the model blindly predicts "1", the accuracy will still be 90\% regardless of the model's actual performance.
        \item \textbf{Define precision, recall, and explain how they combine to form the F1-score.}
        \\ Precision: $\frac{TP}{TP + FP}$, Recall: $\frac{TP}{TP + FN}$, F1-score: $2 \cdot \frac{Precision \cdot Recall}{Precision + Recall}$.
        \\ The F1 score is the harmonic mean of precision and recall that is not biased to either metric.
        \item \textbf{Explain why the macro-average F1-score is a suitable metric for imbalanced classification.}
        \\ Macro average F1 score averages the F1 score for each class.
            This approach gives equal weight to each class, so one class does not dominate the overall score.
            If we simply use accuracy, the model could achieve by simply predicting all classses as 1.
            If macro average F1 score, the model has to also perfrom well on "not 1 class", 
            because the score for "not 1 class" will also be averaged in the final score.
    \end{enumerate}
    \vspace{1em} % Remove this if you don't need extra space
    
    \item \textbf{Missing Values (Report Only)}
    % TODO: 
        \begin{enumerate}[label=(\arabic*)]
        \item \textbf{Choose an imputation strategy (e.g., mean imputation, median imputation, or KNN imputation).}
        \\ KNN imputation.
        \item \textbf{Justify why your chosen strategy would be appropriate for this dataset.}
        \\ Since we are dealing with handwriting dataset(Digits dataset), the spatial relationships becomes important when it comes to missing value interpolation.
            Thus, KNN is the appropriate option as it can capture the features near the missing pixel.
        \item \textbf{Discuss the potential drawbacks of an alternative method.}
        \\ For instance, mean imputation can introduce bias. If we take mean value of the entire dataset and replace it with the pixel that is missing,
            the imputed value may not reflect the actual value of that pixel for a specific sample because the mean shuffles pixel values that are not relevant to the missing pixel.

        \end{enumerate}
            
\end{enumerate}

\end{document}